% GENERATED FILE. Edit canonical lessons, then run npm run generate:books.
% canonical-source-hash: fnv1a64:cb5e71fc40c8c2a5
% canonical-lessons: ZH-C498-tiyu, ZH-C498-qiuchang, ZH-W498-qiu-ball, ZH-C498-qiumen, ZH-C498-shemen, ZH-C498-jinqiu

\chapter{On the Pitch}
\label{ch:zh-on-the-pitch}
\hlchaptermodality{498}

\begin{chapteropening}
\textbf{By the end of this chapter:} \emph{I can talk about PE, a pitch, the ball, the goalmouth and scoring a goal.}

Complete the last lesson of chapter 498: I can talk about PE, a pitch, the ball, the goalmouth and scoring a goal.
\end{chapteropening}

\section[tǐyù]{tǐyù — physical education (a school subject)}

\label{lesson:ZH-C498-tiyu}



\begin{quote}
Before the new one: say the Mandarin for knowledge, then the Mandarin for to go on to further study.
\end{quote}

\subsection*{You'll want to know: tǐyù}
\emph{tǐyù} — \textquotedblleft{}physical education (a school subject)\textquotedblright{}.

\begin{quote}
\emph{wǒ zuì xǐhuan de kè shì tǐyù} — My favourite lesson is PE.
\end{quote}

Its characters come later; for now it is a word for the ear.

\begin{cousinweb}
\emph{tǐyù} is built from \zh{体} 'the body' and \emph{yù} 'to rear', both inherited from Middle Chinese, \emph{thejX} and \emph{yuwk} (Guangyun, in Baxter's transcription), which glosses the first 'the body' and the second 'to nourish' and 'to grow'.
\end{cousinweb}

\begin{grammarlens}[title={What you've built}]
One more word for saying how it was.
\end{grammarlens}

\subsection*{Your turn}
\begin{itemize}
\raggedright
  \item \emph{Say it:} \emph{tǐyù}
  \item \emph{Say it:} \emph{tǐyù}, once more
  \item \emph{Say it:} \emph{wǒ zuì xǐhuan de kè shì tǐyù} — My favourite lesson is PE.
  \item \emph{Recall:} say the Mandarin for knowledge, then the Mandarin for to go on to further study, then say \emph{tǐyù} again
\end{itemize}

\begin{tcolorbox}[breakable,colback=teal!4,colframe=teal!35!black,title={Before you move on}]
What does tǐyù mean? (\textquotedblleft{}Physical education (a school subject)\textquotedblright{}.) Say it once more.
\end{tcolorbox}
\section[\texorpdfstring{\zh{球场} (qiúchǎng)}{球场 (qiúchǎng)}]{\zh{球场} (qiúchǎng) — sports ground; pitch; court}

\label{lesson:ZH-C498-qiuchang}



\begin{quote}
Before the new one: say the Mandarin for to go on to further study, then the Mandarin for physical education.
\end{quote}

\subsection*{You'll want to know: \zh{球场}}
\textbf{\zh{球场}} — \emph{qiúchǎng} — \textquotedblleft{}sports ground; pitch; court\textquotedblright{}.

\begin{quote}
\emph{wǒmen zài qiúchǎng tī le yí ge xiàwǔ zúqiú} — We played football on the pitch all afternoon.
\end{quote}

The new character, in pieces: \zh{球} — eleven strokes in two pieces: on the left the jade piece of \zh{理} \emph{lǐ}, four strokes; on the right seven strokes, a bar, a long vertical with a hook, a dot and a rising stroke on the left, a short and a long falling stroke on the right, and a last dot at the top right. The next lesson puts it in your hand.

\begin{cousinweb}
\textbf{\zh{球场}} (\emph{qiúchǎng}) is built from \emph{qiú} 'a ball' and \zh{场} 'a ground', both inherited from Middle Chinese, \emph{gjuw} and \emph{drjang} (Guangyun, in Baxter's transcription), which glosses the second 'a threshing floor'. It glosses the first character 'a fine jade'; the sense 'a ball' belongs to a second character with the same sound, glossed 'a ball of hair to be struck', which modern writing spells the same way.
\end{cousinweb}

\begin{grammarlens}[title={What you've built}]
One more word for saying how it was.
\end{grammarlens}

\subsection*{Your turn}
\begin{itemize}
\raggedright
  \item \emph{Say it:} \emph{qiúchǎng}
  \item \emph{Say it:} \emph{qiúchǎng}, once more
  \item \emph{Say it:} \emph{wǒmen zài qiúchǎng tī le yí ge xiàwǔ zúqiú} — We played football on the pitch all afternoon.
  \item \emph{Recall:} say the Mandarin for to go on to further study, then the Mandarin for physical education, then say \emph{qiúchǎng} again
\end{itemize}

\begin{tcolorbox}[breakable,colback=teal!4,colframe=teal!35!black,title={Before you move on}]
What does \zh{球场} mean? (\textquotedblleft{}Sports ground; pitch; court\textquotedblright{}.) Say it once more.
\end{tcolorbox}
\section[\texorpdfstring{\zh{球} (qiú)}{球 (qiú)}]{\zh{球} — the jade piece, then seven strokes}

\label{lesson:ZH-W498-qiu-ball}



\begin{quote}
The word first, then the character inside it.

\begin{itemize}
\raggedright
  \item \emph{Recall:} say \emph{qiúchǎng} — \textbf{R1}, one lesson back, before it decays
\end{itemize}

\begin{quote}
\zh{理} \zh{场}
\end{quote}

\emph{Read it:} each one aloud, with nothing said first — \emph{lǐ}, \emph{chǎng}
\end{quote}

\subsection*{Script — the jade piece, then seven strokes}
\hlblockfigure{\detokenize{figures/ZH-W498-qiu-ball-filmstrip.pdf}}{How \zh{球} is written, stroke by stroke}

\begin{quote}
\zh{球}
\end{quote}

\textbf{Two pieces, left before right.}

\textbf{The left piece, four strokes}, as on the left of \zh{理}: two bars, a vertical through them, then a stroke rising to the right.

\textbf{The right piece, seven strokes.} A bar. A long vertical down through it that hooks at the foot. On the left, a dot and a stroke rising to the right. On the right, a short stroke falling to the left and a long stroke falling to the right. Last, the small dot at the top right.

4 + 7 = eleven.

\subsection*{Your turn}
\begin{itemize}
\raggedright
  \item \emph{Write it:} \zh{球}, all eleven, with the model in front of you
  \item \emph{Cover:} write it again from memory, and count the pieces, not the strokes
\end{itemize}

\begin{tcolorbox}[breakable,colback=teal!4,colframe=teal!35!black,title={Before you move on}]
How many pieces does \zh{球} have? (\textbf{Two}.) Which character you know has the same left piece? (\textbf{\zh{理}}.) Which stroke comes last? (\textbf{The small dot at the top right}.)

Source: \href{https://raw.githubusercontent.com/chanind/hanzi-writer-data/68d10a4b21150cae5e1ebbd223eed289cf32d90c/data/\%E7\%90\%83.json}{Hanzi Writer Data \zh{球}, pinned stroke paths and medians}.
\end{tcolorbox}
\section[\texorpdfstring{\zh{球门} (qiúmén)}{球门 (qiúmén)}]{\zh{球门} (qiúmén) — goalmouth; goal (in football)}

\label{lesson:ZH-C498-qiumen}



\begin{quote}
Before the new one: say the Mandarin for physical education, then the Mandarin for a sports pitch.
\end{quote}

\subsection*{You'll want to know: \zh{球门}}
\textbf{\zh{球门}} — \emph{qiúmén} — \textquotedblleft{}goalmouth; goal (in football)\textquotedblright{}.

\begin{quote}
\textbf{\zh{他站在球门前边}} \emph{tā zhàn zài qiúmén qiánbian} — He's standing in front of the goal.
\end{quote}

\begin{cousinweb}
\textbf{\zh{球门}} (\emph{qiúmén}) is built from \zh{球} 'a ball' and \zh{门} 'a gate', both inherited from Middle Chinese, \emph{gjuw} and \emph{mwon} (Guangyun, in Baxter's transcription), which says the character is made of two door leaves.
\end{cousinweb}

\begin{grammarlens}[title={What you've built}]
One more word for saying how it was.
\end{grammarlens}

\subsection*{Your turn}
\begin{itemize}
\raggedright
  \item \emph{Say it:} \emph{qiúmén}
  \item \emph{Say it:} \emph{qiúmén}, once more
  \item \emph{Say it:} \emph{tā zhàn zài qiúmén qiánbian} — He's standing in front of the goal.
  \item \emph{Recall:} say the Mandarin for physical education, then the Mandarin for a sports pitch, then say \emph{qiúmén} again
\end{itemize}

\begin{tcolorbox}[breakable,colback=teal!4,colframe=teal!35!black,title={Before you move on}]
What does \zh{球门} mean? (\textquotedblleft{}Goalmouth; goal (in football)\textquotedblright{}.) Say it once more.
\end{tcolorbox}
\section[shèmén]{shèmén — to shoot at goal}

\label{lesson:ZH-C498-shemen}



\begin{quote}
Before the new one: say the Mandarin for a sports pitch, then the Mandarin for the goalmouth.
\end{quote}

\subsection*{You'll want to know: shèmén}
\emph{shèmén} — \textquotedblleft{}to shoot at goal\textquotedblright{}.

\begin{quote}
\emph{tā shèmén de shíhou qiúmén qiánbian méiyǒu rén} — When he shot, there was nobody in front of the goal.
\end{quote}

Its characters come later; for now it is a word for the ear.

\begin{cousinweb}
\emph{shèmén} is built from \emph{shè} 'to shoot' and \zh{门} 'a gate', both inherited from Middle Chinese, \emph{zyaeH} and \emph{mwon} (Guangyun, in Baxter's transcription), which glosses the first with the arts of archery.
\end{cousinweb}

\begin{grammarlens}[title={What you've built}]
One more word for saying how it was.
\end{grammarlens}

\subsection*{Your turn}
\begin{itemize}
\raggedright
  \item \emph{Say it:} \emph{shèmén}
  \item \emph{Say it:} \emph{shèmén}, once more
  \item \emph{Say it:} \emph{tā shèmén de shíhou qiúmén qiánbian méiyǒu rén} — When he shot, there was nobody in front of the goal.
  \item \emph{Recall:} say the Mandarin for a sports pitch, then the Mandarin for the goalmouth, then say \emph{shèmén} again
\end{itemize}

\begin{tcolorbox}[breakable,colback=teal!4,colframe=teal!35!black,title={Before you move on}]
What does shèmén mean? (\textquotedblleft{}To shoot at goal\textquotedblright{}.) Say it once more.
\end{tcolorbox}
\section[\texorpdfstring{\zh{进球} (jìnqiú)}{进球 (jìnqiú)}]{\zh{进球} (jìnqiú) — to score a goal}

\label{lesson:ZH-C498-jinqiu}



\begin{quote}
Before the new one: say the Mandarin for the goalmouth, then the Mandarin for to shoot at goal.
\end{quote}

\subsection*{You'll want to know: \zh{进球}}
\textbf{\zh{进球}} — \emph{jìnqiú} — \textquotedblleft{}to score a goal\textquotedblright{}.

\begin{quote}
\textbf{\zh{进球的时候我感觉很开心}} \emph{jìnqiú de shíhou wǒ gǎnjué hěn kāixīn} — When I score, I feel great.
\end{quote}

\begin{cousinweb}
\textbf{\zh{进球}} (\emph{jìnqiú}) is built from \zh{进} 'to go in' and \zh{球} 'a ball', both inherited from Middle Chinese, \emph{tsinH} and \emph{gjuw} (Guangyun, in Baxter's transcription), which glosses the first 'forward'.
\end{cousinweb}

\begin{grammarlens}[title={What you've built}]
One more word for saying how it was.
\end{grammarlens}

\subsection*{Your turn}
\begin{itemize}
\raggedright
  \item \emph{Say it:} \emph{jìnqiú}
  \item \emph{Say it:} \emph{jìnqiú}, once more
  \item \emph{Say it:} \emph{jìnqiú de shíhou wǒ gǎnjué hěn kāixīn} — When I score, I feel great.
  \item \emph{Recall:} say the Mandarin for the goalmouth, then the Mandarin for to shoot at goal, then say \emph{jìnqiú} again
\end{itemize}

\begin{tcolorbox}[breakable,colback=teal!4,colframe=teal!35!black,title={Before you move on}]
What does \zh{进球} mean? (\textquotedblleft{}To score a goal\textquotedblright{}.) Say it once more.
\end{tcolorbox}
